\chapter*{Abstract}
    \addcontentsline{toc}{chapter}{Abstract}
A vision-language model (VLM) is able to read a chart, document, scanned form, or natural image and answer questions about it. 
But how do we know if the answer is actually derived from the image itself, rather than being a lucky guess on a question based on the model's remembered language patterns?
The correctness of the answer itself is not enough if the reasoning for that answer is not grounded in the visual evidence of the image.
Charts and documents are a good example of this problem, because to answer the question correctly, the model has to actually ground its reasoning on specific parts of the image,
making it easier to spot when the model is looking at the wrong thing. Although recent work has made progress on this issue,
there was no isolation of the training objective as the main variable, with all other factors such as the model, the data, and the budget held constant.
Contrastive, generative, and explanation-based methods have all been proposed, but they have not been directly compared under the same conditions,
so it is not clear what the objective itself contributed to the final performance. 
In this paper, we keep the pipeline fixed and change only the training signal, so we can see what the objective does to answer accuracy, explanation quality, 
and evidence sensitivity when the model, the data, and the compute budget are all held steady.
We compare five settings head to head rather than treating them as unrelated systems: BLIP-2 generative fine-tuning, BLIP-2 contrastive-enhanced fine-tuning,
Qwen2-VL under answer-only, rationale-generative, and explanation-aware supervision.
All the runs use the same resource-constrained single-seed protocol, and the same evaluation code as much as possible,
across ScienceQA, A-OKVQA, ChartQA, DocVQA, and a VQAv2 subset.
The picture that emerges is mixed. On ScienceQA, Qwen2-VL-2B reaches 0.800 accuracy with answer-only supervision,
0.780 in both rationale-generative and length-aware explanation control settings, and 0.765 with fixed-alpha explanation-aware training.
On A-OKVQA, the same model reaches 0.830 with answer-only supervision, 0.790 in both rationale-generative and length-aware explanation control training,
and 0.795 with fixed-alpha explanation-aware training.
ChartQA, DocVQA, and the VQAv2 subset have no gold rationales, so those runs fall back to answer-only controls;
they lift the native headline metrics, but they do not test rationale supervision at all.
BLIP-2 barely moves: the contrastive objective gives a small gain compared to the generative one on ScienceQA 0.480 vs 0.475,
but the retrieval diagnostic stays near frozen and random controls in both cases.
Scale is the exception. With a fixed explanation-aware objective on ScienceQA, Qwen2-VL-7B reaches 0.885 while the 2B version reaches 0.765.
The masking analysis tells us whether the model relies on the image, not whether its explanation is faithful.
Overall, the results do not show a clear win for explanation-aware training.
The rationale supervision can help on a single-GPU budget, but the gain is easy to mis-credit to scale or tuning.
Hence, it only stands up when answer-only controls, objective balance, model scale, transfer and evidence sensitivity are each reported separately.

    \clearpage
